\properlane{mercy-received}{The forgiven servant must become merciful}

A man who has received mercy can still live as though mercy had never
been shown him. That is the Gospel's disturbing possibility. The first
servant's release is neither imaginary nor earned by his conduct toward
the second. It precedes that conduct. The question is what the recipient
does with a freedom he could not have purchased. The Mass brings that
question into a household asking protection, hearing of spiritual warfare,
offering sacrifice, and receiving immortal nourishment. God's gifts are
first; their recipient's response nevertheless matters.

\subsection{Why the enormous debt is more than an impressive number}
The parable begins with the king, not with the servant's grievance.
That order prevents the hearer from measuring everything by what another
person owes him. The first debt is immense, and the text expressly says
that the servant has no means to pay. He promises everything, but his
promise does not produce the missing resources. The king's compassion
does what further time could not do: it cancels the debt. The man leaves
because he has been released.

Chrysostom makes this disproportion theological by recounting gifts that
precede any human repayment. God created us, gave the world in which we
live, and, after human ingratitude, sent his Son, opened heaven, gave
baptismal remission, and bestowed the Spirit. He is not calculating an
exchange rate between talents and present-day income. He is asking what
a creature and redeemed person could possibly offer that would put God
in his debt. Even a life of grateful service remains something by which
God benefits the servant.\footnote{Chrysostom, \work{Matthew} LXI,
\S\S1--2, NPNF I.10.}

This explains why the Introit's universal creator belongs at the
beginning of a reading centered on pardon. The petitioner stands before
the Lord of all, not an equal negotiating with another proprietor.
Mordecai invokes that lordship in danger; the Church receives his prayer
without becoming owner of the God she invokes. The Sunday Collect
similarly calls the assembly God's family and asks continuing mercy.
The family does not first establish its independence and then purchase
occasional help. Its very life of good works unfolds under protection
it continues to need.

Chrysostom also explains the initial threat of sale. He reads it as a
means of bringing the servant to supplication and awareness of his
liability. His evidence is the outcome: the king does not execute the
sale, but grants more than the petition sought. The threatened loss
teaches the size of the rescue. Thus the first scene gives the servant
both freedom and a concrete memory with which to recognize his
neighbor's need.\footnote{Chrysostom, \work{Matthew} LXI, \S3.}

\subsection{The words he should have recognized}
The second servant repeats the first servant's appeal: patience, with
a promise of payment. The gestures echo as well; the debtor falls down
and begs. Chrysostom lingers over that recurrence. The forgiven man
should recognize the harbor through which he has just escaped. Instead
he seizes a fellow servant by the throat. His demand sounds like justice
when isolated---pay what you owe---but its use inside this story reveals
a heart refusing the pattern of the mercy it has received.

The comparison does not say the lesser debt is nonexistent. Nor does
it say that a wound inflicted by a neighbor cannot be grave. It makes
the two servants equals before their master and places their encounter
after an immeasurable gift. The second servant asks only for delay;
the first refuses even that. Chrysostom's inference depends upon the
asymmetry: complete forgiveness received, minimal patience denied.
The prison is therefore more than an unfortunate sequel. It exposes
the recipient's choice to use his new freedom against another.\footnote{Chrysostom, \work{Matthew} LXI, \S4.}

Augustine brings this encounter into the hearer's daily prayer. Every
person is simultaneously God's debtor and someone else's creditor.
Neither identity can be reserved for a separate compartment of life.
The person who asks God for pardon while refusing all pardon to a
neighbor has made two contradictory appeals. Augustine also calls
the petitioner God's beggar: one person asks us for bread, while we
ask God for Christ, the living bread. Giving and forgiving belong to
a creature who must continue to receive.\footnote{Augustine,
Sermon~83 \S\S1--4, NPNF I.6, Sermon XXXIII.}

This is a substantial addition to Chrysostom's comparison of the two
scenes. Chrysostom teaches the servant to remember a decisive gift;
Augustine teaches the community to recognize its present dependence
each time it prays. The Gradual's refuge ``from generation to
generation'' gives that dependence its enduring home. The eternal God
does not expire when our resolve does. And the Introit's undefiled
walk becomes a petition for grace-enabled conduct, as Augustine's
Aleph exposition teaches, rather than a platform from which the
apparently righteous can despise other debtors.

\subsection{Why the pardon was real, and why the ending remains fearful}
The king says, ``I forgave thee all the debt.'' Chrysostom insists upon
the reality of the gift and speaks sharply of its reversal: the servant
himself has revoked the sentence by his cruelty. Augustine likewise
warns the baptized that the forgiven debt can be exacted again. Their
language makes the warning severe. It does not describe an initial
pardon that was only a performance concealing a plan never to forgive.\footnote{Chrysostom, \work{Matthew} LXI, \S4; Augustine,
Sermon~83 \S\S6--7.}

The ending's ``until'' is equally severe. In an ordinary loan, until
payment might designate the date on which a prisoner goes free. Here
the impossible liability has already been established. Chrysostom
therefore reads the sentence as unending, not a promise that the
servant will finally produce the money. The force of Jesus' conclusion
is addressed to the disciples themselves: they must forgive from the
heart. Treating the last scene as an optional exaggeration would remove
the very point that the king's question presses.

Aquinas later asks with greater technical precision how sins already
forgiven can return. He refuses the idea that a human act can undo
God's completed act of pardon. He also refuses an account in which
remission was only conditional and never truly effected by grace.
A new mortal sin incurs a new debt of eternal punishment. Previous
gifts matter because contempt of greater mercy aggravates the new
offense; the old sinful acts do not simply become unforgiven again.\footnote{Thomas Aquinas, \work{Summa theologiae} III, q.~88,
a.~1, answer and replies, Fathers of the English Dominican Province,
second revised edition (1920).}

Article~3 adds a further distinction. The new liability is not
automatically the numerical sum of every old penalty plus the new
one. Its measure follows the new sin's gravity, with the received
favor contributing to the gravity of ingratitude. Article~2 explains
why hatred of one's brother is especially pointed: it contradicts
the form of the favor received, the remission of what was owed.
Aquinas thus explains the renewed accountability without treating
forgiveness as unreal. His distinctions are a later doctrinal account;
Chrysostom and Augustine themselves preach the warning in the stronger
images of the parable.\footnote{Aquinas, III, q.~88, aa.~2--4,
especially a.~2, answer, and a.~3, answer and reply to objection~3.}

The Sunday Secret holds the same priority of gift in sacramental
prayer. God has willed the offering's propitiatory and restoring
effects. Mercy toward a neighbor does not buy that divine initiative.
Rather, the recipient is called to live truthfully within it. The
Trinity Preface prevents the sequence from being imagined as a
conflict within God: the Father, Son, and Spirit are confessed in
equal majesty. The king's mercy and judgment cannot be assigned to
two competing gods.

\subsection{A brother is not the enemy to be destroyed}
Ephesians identifies the conflict as something other than wrestling
against flesh and blood. Chrysostom's exposition of the preceding
verse~9 explicitly recalls the unmerciful servant when it tells
masters that they too have a Master. His warning is not merely that
everyone should be more polite. Rank cannot exempt a person from
the common Lord's judgment. A servant's injury is not morally
invisible because he has less earthly power.\footnote{Chrysostom,
\work{Ephesians} XXII, on 6:9.}

Chrysostom then redirects anger toward the spiritual enemy. The
weapons are truth, justice, faith, salvation, and peace with God.
Revenge on the fellow servant would imitate the enemy's destructive
purpose while pretending to win a victory. Forgiveness is therefore
part of the spiritual battle, not surrender to it. One can stand
firmly against falsehood without nourishing a wish for the liar's
ruin. One can oppose an unjust act without turning its perpetrator
into a being beyond the claims of charity.

Augustine's final paragraph on the Gospel protects this point from
another mistake. He distinguishes inward pardon from abandonment
of correction. Love can require rebuke directed toward the person's
good. His examples reflect ancient disciplinary practices; the
distinction does not authorize their indiscriminate imitation. It
does establish that forgiveness cannot simply mean declaring every
act harmless or refusing to prevent its repetition. Truth, restitution,
protection, and proportionate accountability remain compatible with
relinquishing hatred.\footnote{Augustine, Sermon~83 \S8. The applications
to restitution and protective boundaries here develop that distinction.}

Job supplies the indispensable counterexample to a cruel use of this
Gospel. Not every sufferer is the condemned servant. The Offertory
names an upright man before it names his wounds; Gregory explains
why they need not punish an antecedent offense. The afflicted person
must not be told that pain proves unforgiveness or a hidden crime.
Nor does involuntary distress reproduce the servant's deliberate
``he would not.'' The Gospel judges a refusal of mercy; it is not
a diagnosis of the speed with which an injured person's emotions heal.

\subsection{Mercy can ask for judgment and still receive nourishment}
At Communion the people ask when God will judge their persecutors.
That question has not been erased by the Gospel. Augustine's martyrs
await God's judgment while remaining faithful to his truth. Their
appeal differs from the servant's seizure of another's throat:
judgment is entrusted to God, and injustice is named rather than
perpetuated. Forgiveness does not require calling persecution good.
The Church can hope for an end to it without making hatred her food.

The Sunday Postcommunion asks precisely for an interior fruit of
what has been received. The mouth has received; the mind must follow.
Schuster describes the sacramental bond reaching mind, heart, and
will. Augustine's distinction between a soft voice and a relentless
heart now acquires a liturgical application. The required transformation
is neither mere verbal politeness nor confidence in an outward act
detached from the person who receives it.

The Alleluia's Exodus strengthens this reading. Augustine understands
the sea's passage as genuine baptismal deliverance and also recalls
the later wilderness falls. Real grace begins a life of fidelity.
Likewise the pardoned servant's departure begins a new obligation;
it does not conclude his moral history. The food of immortality is
received along that history, strengthening the family toward a
destination it has not yet reached.

When the mission prayers are appointed, mercy's outward movement
becomes explicit. The Collect asks workers to carry the word to all
peoples; the Secret remembers redemption for all and worship across
the earth; the Postcommunion asks continuing growth in true faith.
A family that receives such a gift cannot regard salvation as its
private possession. The human nations are recipients of proclamation,
not targets of the Epistle's spiritual weapons.

Where Luke is commemorated, his Collect names costly service for
God's name. The Secret's free mind is freedom for service rather
than freedom to imprison another; its healing and glory are gifts
sought through the offering. The Postcommunion asks the altar's
gift to sanctify the soul and keep it safe. Those prayers extend
the Sunday's inward demand into the communion of saints: received
life becomes service, and service remains dependent upon grace.

\begin{fourSenses}
\item[Literal.] Mordecai petitions for endangered Israel; Paul equips a
community for spiritual resistance; the king really cancels an unpayable
debt before the servant refuses a fellow debtor; upright Job really
suffers. Their different situations remain distinct.
\item[Allegorical.] Christ frees and gathers his household. Augustine's
baptismal Exodus and his account of the pardoned debtor illuminate a
people released by grace and fed with the living bread. Schuster's
Job-Christ typology places innocent suffering beside the restoring
sacrifice.
\item[Moral.] Remember received pardon, forgive inwardly, oppose falsehood
and cruelty without hatred, and let sacramental reception become merciful
conduct. Correction and protection need not be abandoned.
\item[Anagogical.] The final accounting exposes the gravity of refusing
mercy. The persecuted await God's judgment, and immortal nourishment
orders the household toward lasting communion with him.
\end{fourSenses}

\clearpage
